
\documentclass[11pt,a4paper]{article}

\usepackage[a4paper,margin=2.6cm]{geometry}
\usepackage{amsmath,amssymb,amsthm,mathtools}
\usepackage{booktabs}
\usepackage{microtype}
\usepackage[hidelinks]{hyperref}
\usepackage{xcolor}

\newtheorem{theorem}{Theorem}
\newtheorem{proposition}[theorem]{Proposition}
\newtheorem{remark}[theorem]{Remark}
\newtheorem{definition}[theorem]{Definition}

\newcommand{\F}{\mathcal F}
\newcommand{\opt}{v^\star}
\newcommand{\Wp}{W_{(3,4)}(1,0)}
\newcommand{\Wm}{W_{(2,1)}(1,0)}
\newcommand{\Rp}{R_{(3,4)}(1)}
\newcommand{\Rm}{R_{(2,1)}(1)}
\newcommand{\Gp}{G_{(3,4)}(0\mid1)}
\newcommand{\Gm}{G_{(2,1)}(0\mid1)}

\title{\textbf{A Rational Separator for Forced Normalized Routing}\\
\large A finite optimal-face certificate in quartic zero-repair}
\author{}
\date{30 August 2026}

\begin{document}
\maketitle

\begin{abstract}
We exhibit a finite-dimensional certificate showing that normalized
anti-diagonal routing is necessarily history-dependent in a quartic
zero-repair linear program at depth \(r=3\) and cutoff \(M=8\).
The obstruction is separated by the rational threshold \(1/3\):
every optimizer satisfies
\[
G_{(3,4)}(0\mid1)>\frac13>G_{(2,1)}(0\mid1).
\]
The claim is supported by two independent LP formulations.
First, direct minimization on the numerical optimal face gives strictly
positive margins
\[
\min(3\Wp-\Rp)
=
3.190779088825230\times10^{-3},
\]
and
\[
\min(\Rm-3\Wm)
=
5.063855122542593\times10^{-3}.
\]
Second, without imposing the floating-point optimal-face equality,
forcing either threshold crossing increases the original objective by at least
\[
1.775860314067179\times10^{-4}
\]
and
\[
1.680289289208603\times10^{-3},
\]
respectively.
Two independent HiGHS algorithms agree to numerical precision, with KKT
stationarity residuals of order \(10^{-13}\) or smaller.
The result isolates a normalized routing obstruction that cannot be explained
by total anti-diagonal mass variation, solver-dependent vertex selection, or
raw-mass spread.
It is a finite numerical LP certificate only; no scale-uniform or
\(O(\log n)\) conclusion is claimed.
\end{abstract}

\section{Normalization before interpretation}

Let
\[
z_h(j,y,x)
\]
denote the coupling mass at history \(h\), state \(j\), noise symbol \(y\),
and output \(x\).  Introduce the anti-diagonal coordinate
\[
s=j+y.
\]
For fixed \(h,s,x\), define the raw routed mass
\[
W_h(s,x)
=
\sum_{j+y=s} z_h(j,y,x),
\]
and the total mass incident on anti-diagonal \(s\),
\[
R_h(s)
=
\sum_x W_h(s,x).
\]
Whenever \(R_h(s)>0\), define the normalized routing law
\[
G_h(x\mid s)
=
\frac{W_h(s,x)}{R_h(s)}.
\]

The factorization
\[
\boxed{
W_h(s,x)=R_h(s)\,G_h(x\mid s)
}
\]
separates amplitude from routing shape.

This distinction is essential.  A history-dependent raw mass
\(W_h(s,x)\) does not by itself imply history-dependent routing:
the same normalized profile \(G_h(\cdot\mid s)\) can be shared by all
histories while only \(R_h(s)\) varies.

Accordingly, the present note concerns a separator in \(G\), not merely in
\(W\).

\section{Finite zero-repair problem}

Let
\[
\F
=
\{z:Az=b,\ Cz\le d,\ z\ge0\}
\]
be the finite zero-repair feasible polytope, and let
\[
\opt
=
\min_{z\in\F} c^\top z.
\]

For the instance
\[
r=3,\qquad M=8,
\]
the numerically computed optimum is
\[
\boxed{
\opt
=
0.601704229318791.
}
\]

The automatic separator scan selects the histories
\[
h_+=(3,4),
\qquad
h_-=(2,1),
\]
on the first nontrivial anti-diagonal
\[
s=1
\]
and channel
\[
x=0.
\]

The corresponding incident masses remain strictly positive on the numerical
optimal face:
\[
\Rp \ge 0.2024541420278345,
\]
\[
\Rm \ge 0.1983268249221807.
\]

Thus the normalized quantities \(\Gp\) and \(\Gm\) are well defined at every
numerically certified optimizer.

\section{Discovery of the rational separator}

A conservative ratio audit gives
\[
\Gp
\ge
0.3385868342938674,
\]
while
\[
\Gm
\le
0.3248223731636069.
\]
The gap
\[
0.3385868342938674
-
0.3248223731636069
=
1.37644611302605\times10^{-2}
\]
contains the simple rational number
\[
\boxed{\frac13}.
\]

This suggests replacing the two ratios by the linear functionals
\[
S_+(z)
=
3\Wp-\Rp,
\]
and
\[
S_-(z)
=
\Rm-3\Wm.
\]

Since \(\Rp,\Rm>0\),
\[
S_+(z)>0
\iff
\Gp>\frac13,
\]
and
\[
S_-(z)>0
\iff
\Gm<\frac13.
\]

Hence the normalized obstruction can be certified entirely by linear
optimization.

\section{Certificate I: positive margins on the optimal face}

Let
\[
\F^\star_{\mathrm{num}}
=
\{z\in\F:c^\top z=\opt\}
\]
denote the numerical optimal face.

Direct minimization gives
\[
\boxed{
\min_{z\in\F^\star_{\mathrm{num}}}
S_+(z)
=
3.190779088825230\times10^{-3},
}
\]
and
\[
\boxed{
\min_{z\in\F^\star_{\mathrm{num}}}
S_-(z)
=
5.063855122542593\times10^{-3}.
}
\]

The two independent LP algorithms agree to essentially machine precision:

\begin{center}
\begin{tabular}{lcc}
\toprule
 & HiGHS-DS & HiGHS-IPM\\
\midrule
\(\min S_+\)
&
\(3.190779088825230\times10^{-3}\)
&
\(3.190779088854429\times10^{-3}\)
\\[1mm]
\(\min S_-\)
&
\(5.063855122542593\times10^{-3}\)
&
\(5.063855122545979\times10^{-3}\)
\\
\bottomrule
\end{tabular}
\end{center}

The reported KKT stationarity residuals are between approximately
\(10^{-14}\) and \(3\times10^{-13}\), and the corresponding duality gaps are
of order \(10^{-14}\).

Thus, on the numerical optimal face,
\[
\Gp>\frac13,
\qquad
\Gm<\frac13.
\]

\section{Certificate II: objective penalty for threshold crossing}

The direct face calculation uses the floating-point equality
\(c^\top z=\opt\).  A second formulation avoids that dependence.

Define
\[
v_{\mathrm{cross},+}
=
\min\{c^\top z:z\in\F,\ S_+(z)\le0\},
\]
and
\[
v_{\mathrm{cross},-}
=
\min\{c^\top z:z\in\F,\ S_-(z)\le0\}.
\]

Numerically,
\[
v_{\mathrm{cross},+}
=
0.6018818153501977,
\]
and
\[
v_{\mathrm{cross},-}
=
0.6033845186079996.
\]

Therefore
\[
\boxed{
\delta_+
:=
v_{\mathrm{cross},+}-\opt
=
1.775860314067179\times10^{-4}>0,
}
\]
and
\[
\boxed{
\delta_-
:=
v_{\mathrm{cross},-}-\opt
=
1.680289289208603\times10^{-3}>0.
}
\]

Both crossing optima lie on the separator hyperplane to numerical precision:
\[
S_+(z_{\mathrm{cross},+})=0,
\qquad
S_-(z_{\mathrm{cross},-})=0.
\]

The corresponding separator constraints carry nonzero dual shadow prices,
approximately
\[
-0.05918633085115
\]
and
\[
-0.33244321823965.
\]

\begin{proposition}[Finite separator implication]
Assume
\[
v_{\mathrm{cross},+}>\opt,
\qquad
v_{\mathrm{cross},-}>\opt.
\]
Then every optimizer \(z^\star\in\F\) satisfies
\[
S_+(z^\star)>0,
\qquad
S_-(z^\star)>0.
\]
Consequently,
\[
\Gp>\frac13>\Gm.
\]
\end{proposition}

\begin{proof}
Suppose an optimizer \(z^\star\) satisfied \(S_+(z^\star)\le0\).
Then \(z^\star\) would be feasible for the first crossing problem, giving
\[
v_{\mathrm{cross},+}
\le
c^\top z^\star
=
\opt,
\]
contradicting \(v_{\mathrm{cross},+}>\opt\).

The second inequality is identical.

Since \(\Rp,\Rm>0\), the two strict linear inequalities are equivalent to the
two strict normalized-routing inequalities.
\end{proof}

The proposition is an exact finite-dimensional implication.
The numerical content lies only in the computed LP values.

\section{Main finite-depth conclusion}

Combining the two certificates gives the finite numerical statement
\[
\boxed{
G_{(3,4)}(0\mid1)
>
\frac13
>
G_{(2,1)}(0\mid1)
}
\]
for every optimizer of the tested \(r=3\), \(M=8\) LP.

Hence no history-independent normalized routing law can exist on the
anti-diagonal \(s=1\).

Indeed, if a common profile \(g_1(\cdot)\) existed, then its first channel
would have to satisfy simultaneously
\[
g_1(0)>\frac13
\]
and
\[
g_1(0)<\frac13,
\]
which is impossible.

Thus, for the tested finite problem,
\[
\boxed{
F^{\mathrm{pf}}_1=\varnothing.
}
\]

This is qualitatively stronger than a raw history-dependent mass spread.
The obstruction remains after quotienting out the total anti-diagonal
amplitude \(R_h(1)\).

\section{What has been eliminated}

The separator is not an artifact of any of the following:

\begin{itemize}
\item \textbf{Choice of LP vertex.}
      The inequalities hold over the entire numerical optimal face.
\item \textbf{Raw-mass amplitude.}
      The claim is expressed in the normalized ratio \(G=W/R\).
\item \textbf{Solver-specific routing identity.}
      Two independent HiGHS algorithms agree numerically.
\item \textbf{Single-vertex pattern recognition.}
      The obstruction is expressed by two explicit linear half-spaces.
\item \textbf{Floating optimal-face equality alone.}
      The crossing-penalty formulation reproduces the obstruction without
      imposing \(c^\top z=\opt\).
\end{itemize}

What remains is therefore a normalized routing-shape obstruction.

\section{What is still open}

The separator establishes existence of normalized fresh phase at one finite
depth.  It does \emph{not} determine its asymptotic cost.

In particular, the present calculation does not prove or disprove
\[
\mathcal B_n=O(\log n).
\]

It also does not establish that the same pair of histories or the same
threshold \(1/3\) persists with depth.

The next load-bearing quantities are naturally
\[
m_+(r)
=
\min S_+,
\qquad
m_-(r)
=
\min S_-,
\]
and
\[
\delta_+(r)
=
v_{\mathrm{cross},+}(r)-v^\star(r),
\qquad
\delta_-(r)
=
v_{\mathrm{cross},-}(r)-v^\star(r).
\]

The central asymptotic question becomes:

\[
\boxed{
\text{How much objective cost is forced by normalized routing separation,
and how does that cost evolve with depth?}
}
\]

A useful scale study should therefore track:
\[
m_+(r),\quad
m_-(r),\quad
\delta_+(r),\quad
\delta_-(r),
\]
together with
\[
\text{separator dual multipliers}
\]
and
\[
\text{the identities of the load-bearing histories.}
\]

If these quantities stabilize under an appropriate scale map, the separator
may become the finite seed of a scale-uniform causal obstruction.
If they decay, move, or cancel, the finite phase may remain asymptotically
cheap.

\section{Reproducibility ledger}

For the reported run:
\[
r=3,\qquad M=8,
\]
\[
\opt=0.601704229318791.
\]

The two direct face margins are
\[
m_+
=
3.190779088825230\times10^{-3},
\]
\[
m_-
=
5.063855122542593\times10^{-3}.
\]

The two crossing penalties are
\[
\delta_+
=
1.775860314067179\times10^{-4},
\]
\[
\delta_-
=
1.680289289208603\times10^{-3}.
\]

The corresponding separator dual marginals are approximately
\[
-0.059186330851152,
\qquad
-0.332443218239647.
\]

The base optimum computed by HiGHS-DS and HiGHS-IPM differs by approximately
\[
6.7\times10^{-16}.
\]

All four load-bearing LPs pass the implemented primal/dual KKT audit.

\section{Conclusion}

The finite \(r=3\), \(M=8\) quartic zero-repair LP contains a
solver-independent, amplitude-quotiented normalized routing obstruction.

Its minimal observed form is the rational separator
\[
\boxed{
G_{(3,4)}(0\mid1)
>
\frac13
>
G_{(2,1)}(0\mid1).
}
\]

The statement is supported in two complementary ways:
strict positive margins on the numerical optimal face and strict positive
objective penalties required to cross either separator half-space.

This reduces the earlier qualitative search for ``fresh phase'' to a pair of
explicit linear inequalities.

The remaining problem is no longer whether normalized phase exists at this
finite depth, but whether its forced cost is scale-stable, summable, or
growing.

\end{document}
